\documentclass[tikz,svgnames,border={10 10},convert={density=300, outext=.png}]{standalone}

% 若需要中文，请用 xelatex 编译并取消下面两行注释：
\usepackage{xeCJK}
% \setCJKmainfont{Noto Sans CJK SC}

\usetikzlibrary{positioning,arrows,fit}

\begin{document}
\begin{tikzpicture}[
    box/.style={rectangle,draw,fill=DarkGray!20,rounded corners,
                node distance=1cm,text width=18em,text centered,
                minimum height=2.4em,thick},
    arrow/.style={draw,-latex',thick},
  ]

  % === 主体流程节点（前述模块一致风格，无外框） ===
  \node[box,fill=PeachPuff!60,text width=24em] (input)
    {接收事件与指标输入：内核/交易所事件、性能采样、自定义指标请求};

  \node[box,below=0.6 of input] (eventlog)
    {事件日志记录（log\_event / log.csv 写入缓冲）};

  \node[box,below=0.5 of eventlog] (perfmon)
    {性能指标采集（CPU/内存/IO / 每秒事件数）};

  \node[box,below=0.5 of perfmon] (metric)
    {指标导出与推送（Prometheus 格式 / 告警触发）};

  \node[box,below=0.5 of metric] (crit)
    {需要刷新或存档？};

  % 底部输出
  \node[box,below=1.8 of crit,fill=LightSteelBlue!60,text width=26em] (output)
    {输出：结构化日志文件、监控指标导出、告警消息 / 日志目录归档};

  % === 箭头（直连 + Yes/No 回路） ===
  \path[arrow] (input) -- (eventlog);
  \path[arrow] (eventlog) -- (perfmon);
  \path[arrow] (perfmon) -- (metric);
  \path[arrow] (metric) -- (crit);

  % Yes 分支：直达输出
  \path[arrow] (crit) -- (output)
    node[midway,left=0.1,draw,outer sep=0pt,fill=white] {Yes};

  % No 分支：先竖直向下，再水平右移（No 在水平段下方），最后折返到“事件日志记录”节点
  \path[draw,thick] (crit.south) ++(0,-1em) coordinate(novert);
  \draw[thick] (novert) -- ++(4,0) coordinate(nojump)
        node[midway,below=2pt,draw,outer sep=0pt,fill=white] {No};
  \draw[arrow] (nojump) |- (eventlog.east);

\end{tikzpicture}

\end{document}
